% 按本轮原书第80、81页照片逐项复核；黑字题目，红字答案。
\QHeading{第1题　下面是四年级同学每天读书时间统计表。}
\par\noindent\begin{minipage}[t]{.485\linewidth}
\par\noindent{\centering 四（1）班同学每天读书时间统计表\par}
\par\noindent\begingroup
\begin{tikzpicture}[x=\dimexpr\linewidth-1pt\relax,y=-1mm,draw=black,text=black,line width=.45pt]
\draw (0.000000,0)--(0.000000,17.000);
\draw (0.166667,0)--(0.166667,17.000);
\draw (0.333333,0)--(0.333333,17.000);
\draw (0.500000,0)--(0.500000,17.000);
\draw (0.666667,0)--(0.666667,17.000);
\draw (0.833333,0)--(0.833333,17.000);
\draw (1.000000,0)--(1.000000,17.000);
\draw (0,0.000)--(1,0.000);
\draw (0,10.000)--(1,10.000);
\draw (0,17.000)--(1,17.000);
\node[font=\normalsize,align=center,inner sep=0pt,text width=\dimexpr0.166667\linewidth-2.5mm\relax] at (0.083333,5.000) {时间\\/分};
\node[font=\normalsize,align=center,inner sep=0pt,text width=\dimexpr0.166667\linewidth-2.5mm\relax] at (0.250000,5.000) {合计};
\node[font=\normalsize,align=center,inner sep=0pt,text width=\dimexpr0.166667\linewidth-2.5mm\relax] at (0.416667,5.000) {$31\sim$\\$60$};
\node[font=\normalsize,align=center,inner sep=0pt,text width=\dimexpr0.166667\linewidth-2.5mm\relax] at (0.583333,5.000) {$21\sim$\\$30$};
\node[font=\normalsize,align=center,inner sep=0pt,text width=\dimexpr0.166667\linewidth-2.5mm\relax] at (0.750000,5.000) {$15\sim$\\$20$};
\node[font=\normalsize,align=center,inner sep=0pt,text width=\dimexpr0.166667\linewidth-2.5mm\relax] at (0.916667,5.000) {$15$\\以下};
\node[font=\normalsize,align=center,inner sep=0pt,text width=\dimexpr0.166667\linewidth-2.5mm\relax] at (0.083333,13.500) {人数};
\node[font=\normalsize,align=center,inner sep=0pt,text width=\dimexpr0.166667\linewidth-2.5mm\relax] at (0.250000,13.500) {$30$};
\node[font=\normalsize,align=center,inner sep=0pt,text width=\dimexpr0.166667\linewidth-2.5mm\relax] at (0.416667,13.500) {$9$};
\node[font=\normalsize,align=center,inner sep=0pt,text width=\dimexpr0.166667\linewidth-2.5mm\relax] at (0.583333,13.500) {$15$};
\node[font=\normalsize,align=center,inner sep=0pt,text width=\dimexpr0.166667\linewidth-2.5mm\relax] at (0.750000,13.500) {$5$};
\node[font=\normalsize,align=center,inner sep=0pt,text width=\dimexpr0.166667\linewidth-2.5mm\relax] at (0.916667,13.500) {$1$};
\end{tikzpicture}\par\endgroup

\end{minipage}\hfill\begin{minipage}[t]{.485\linewidth}
\par\noindent{\centering 四（2）班同学每天读书时间统计表\par}
\par\noindent\begingroup
\begin{tikzpicture}[x=\dimexpr\linewidth-1pt\relax,y=-1mm,draw=black,text=black,line width=.45pt]
\draw (0.000000,0)--(0.000000,17.000);
\draw (0.166667,0)--(0.166667,17.000);
\draw (0.333333,0)--(0.333333,17.000);
\draw (0.500000,0)--(0.500000,17.000);
\draw (0.666667,0)--(0.666667,17.000);
\draw (0.833333,0)--(0.833333,17.000);
\draw (1.000000,0)--(1.000000,17.000);
\draw (0,0.000)--(1,0.000);
\draw (0,10.000)--(1,10.000);
\draw (0,17.000)--(1,17.000);
\node[font=\normalsize,align=center,inner sep=0pt,text width=\dimexpr0.166667\linewidth-2.5mm\relax] at (0.083333,5.000) {时间\\/分};
\node[font=\normalsize,align=center,inner sep=0pt,text width=\dimexpr0.166667\linewidth-2.5mm\relax] at (0.250000,5.000) {合计};
\node[font=\normalsize,align=center,inner sep=0pt,text width=\dimexpr0.166667\linewidth-2.5mm\relax] at (0.416667,5.000) {$31\sim$\\$60$};
\node[font=\normalsize,align=center,inner sep=0pt,text width=\dimexpr0.166667\linewidth-2.5mm\relax] at (0.583333,5.000) {$21\sim$\\$30$};
\node[font=\normalsize,align=center,inner sep=0pt,text width=\dimexpr0.166667\linewidth-2.5mm\relax] at (0.750000,5.000) {$15\sim$\\$20$};
\node[font=\normalsize,align=center,inner sep=0pt,text width=\dimexpr0.166667\linewidth-2.5mm\relax] at (0.916667,5.000) {$15$\\以下};
\node[font=\normalsize,align=center,inner sep=0pt,text width=\dimexpr0.166667\linewidth-2.5mm\relax] at (0.083333,13.500) {人数};
\node[font=\normalsize,align=center,inner sep=0pt,text width=\dimexpr0.166667\linewidth-2.5mm\relax] at (0.250000,13.500) {$30$};
\node[font=\normalsize,align=center,inner sep=0pt,text width=\dimexpr0.166667\linewidth-2.5mm\relax] at (0.416667,13.500) {$13$};
\node[font=\normalsize,align=center,inner sep=0pt,text width=\dimexpr0.166667\linewidth-2.5mm\relax] at (0.583333,13.500) {$10$};
\node[font=\normalsize,align=center,inner sep=0pt,text width=\dimexpr0.166667\linewidth-2.5mm\relax] at (0.750000,13.500) {$5$};
\node[font=\normalsize,align=center,inner sep=0pt,text width=\dimexpr0.166667\linewidth-2.5mm\relax] at (0.916667,13.500) {$2$};
\end{tikzpicture}\par\endgroup

\end{minipage}\par
\Problem{将上面的两个表合并成一个表。}
\par\noindent{\centering 四年级同学每天读书时间统计表\par}
\par\noindent\begingroup
\begin{tikzpicture}[x=\dimexpr\linewidth-1pt\relax,y=-1mm,draw=black,text=black,line width=.45pt]
\draw (0.000000,0)--(0.000000,19.500);
\draw (0.166667,0)--(0.166667,19.500);
\draw (0.333333,0)--(0.333333,19.500);
\draw (0.500000,0)--(0.500000,19.500);
\draw (0.666667,0)--(0.666667,19.500);
\draw (0.833333,0)--(0.833333,19.500);
\draw (1.000000,0)--(1.000000,19.500);
\draw (0,0.000)--(1,0.000);
\draw (0,6.500)--(1,6.500);
\draw (0,13.000)--(1,13.000);
\draw (0,19.500)--(1,19.500);
\node[font=\normalsize,align=center,inner sep=0pt,text width=\dimexpr0.166667\linewidth-2.5mm\relax] at (0.083333,3.250) {时间/分};
\node[font=\normalsize,align=center,inner sep=0pt,text width=\dimexpr0.166667\linewidth-2.5mm\relax] at (0.250000,3.250) {合计};
\node[font=\normalsize,align=center,inner sep=0pt,text width=\dimexpr0.166667\linewidth-2.5mm\relax] at (0.416667,3.250) {$31\sim60$};
\node[font=\normalsize,align=center,inner sep=0pt,text width=\dimexpr0.166667\linewidth-2.5mm\relax] at (0.583333,3.250) {$21\sim30$};
\node[font=\normalsize,align=center,inner sep=0pt,text width=\dimexpr0.166667\linewidth-2.5mm\relax] at (0.750000,3.250) {$15\sim20$};
\node[font=\normalsize,align=center,inner sep=0pt,text width=\dimexpr0.166667\linewidth-2.5mm\relax] at (0.916667,3.250) {$15$以下};
\node[font=\normalsize,align=center,inner sep=0pt,text width=\dimexpr0.166667\linewidth-2.5mm\relax] at (0.083333,9.750) {四（1）班};
\node[font=\normalsize,align=center,inner sep=0pt,text width=\dimexpr0.166667\linewidth-2.5mm\relax] at (0.250000,9.750) {{\color{answerred}$30$}};
\node[font=\normalsize,align=center,inner sep=0pt,text width=\dimexpr0.166667\linewidth-2.5mm\relax] at (0.416667,9.750) {{\color{answerred}$9$}};
\node[font=\normalsize,align=center,inner sep=0pt,text width=\dimexpr0.166667\linewidth-2.5mm\relax] at (0.583333,9.750) {{\color{answerred}$15$}};
\node[font=\normalsize,align=center,inner sep=0pt,text width=\dimexpr0.166667\linewidth-2.5mm\relax] at (0.750000,9.750) {{\color{answerred}$5$}};
\node[font=\normalsize,align=center,inner sep=0pt,text width=\dimexpr0.166667\linewidth-2.5mm\relax] at (0.916667,9.750) {{\color{answerred}$1$}};
\node[font=\normalsize,align=center,inner sep=0pt,text width=\dimexpr0.166667\linewidth-2.5mm\relax] at (0.083333,16.250) {四（2）班};
\node[font=\normalsize,align=center,inner sep=0pt,text width=\dimexpr0.166667\linewidth-2.5mm\relax] at (0.250000,16.250) {{\color{answerred}$30$}};
\node[font=\normalsize,align=center,inner sep=0pt,text width=\dimexpr0.166667\linewidth-2.5mm\relax] at (0.416667,16.250) {{\color{answerred}$13$}};
\node[font=\normalsize,align=center,inner sep=0pt,text width=\dimexpr0.166667\linewidth-2.5mm\relax] at (0.583333,16.250) {{\color{answerred}$10$}};
\node[font=\normalsize,align=center,inner sep=0pt,text width=\dimexpr0.166667\linewidth-2.5mm\relax] at (0.750000,16.250) {{\color{answerred}$5$}};
\node[font=\normalsize,align=center,inner sep=0pt,text width=\dimexpr0.166667\linewidth-2.5mm\relax] at (0.916667,16.250) {{\color{answerred}$2$}};
\end{tikzpicture}\par\endgroup

\QHeading{人数核对}



\par\noindent\begin{minipage}[t]{.485\linewidth}
\Answer{四（1）班人数：}
\Answer{\Formula{\begin{aligned}[t]&\phantom{={}}9+15+5+1\\
&=(9+1)+(15+5)\\
&=10+20\\
&=30\text{（人）}\end{aligned}}}

\end{minipage}\hfill\begin{minipage}[t]{.485\linewidth}
\Answer{四（2）班人数：}
\Answer{\Formula{\begin{aligned}[t]&\phantom{={}}13+10+5+2\\
&=(13+2)+(10+5)\\
&=15+15\\
&=30\text{（人）}\end{aligned}}}

\end{minipage}\par
\QHeading{第2题　四年级参加知识竞赛的同学成绩如下。（单位：分）}
\par\noindent\begin{minipage}[t]{.485\linewidth}
\Problem{四（1）班：$85\quad92\quad87\quad87\quad85$\par\hspace*{5.2em}$60\quad95\quad89\quad78\quad72$}

\end{minipage}\hfill\begin{minipage}[t]{.485\linewidth}
\Problem{四（2）班：$66\quad95\quad87\quad98\quad73$\par\hspace*{5.2em}$77\quad85\quad87\quad90\quad92$}

\end{minipage}\par
\Problem{① 根据上面的数据完成统计表。}
\par\noindent{\centering 四年级参加知识竞赛的同学成绩统计表\par}
\par\noindent\begingroup
\begin{tikzpicture}[x=\dimexpr\linewidth-1pt\relax,y=-1mm,draw=black,text=black,line width=.45pt]
\draw (0.000000,0)--(0.000000,19.500);
\draw (0.200000,0)--(0.200000,19.500);
\draw (0.400000,0)--(0.400000,19.500);
\draw (0.600000,0)--(0.600000,19.500);
\draw (0.800000,0)--(0.800000,19.500);
\draw (1.000000,0)--(1.000000,19.500);
\draw (0,0.000)--(1,0.000);
\draw (0,6.500)--(1,6.500);
\draw (0,13.000)--(1,13.000);
\draw (0,19.500)--(1,19.500);
\node[font=\normalsize,align=center,inner sep=0pt,text width=\dimexpr0.200000\linewidth-2.5mm\relax] at (0.100000,3.250) {成绩/分};
\node[font=\normalsize,align=center,inner sep=0pt,text width=\dimexpr0.200000\linewidth-2.5mm\relax] at (0.300000,3.250) {$90$及$90$以上};
\node[font=\normalsize,align=center,inner sep=0pt,text width=\dimexpr0.200000\linewidth-2.5mm\relax] at (0.500000,3.250) {$80\sim89$};
\node[font=\normalsize,align=center,inner sep=0pt,text width=\dimexpr0.200000\linewidth-2.5mm\relax] at (0.700000,3.250) {$70\sim79$};
\node[font=\normalsize,align=center,inner sep=0pt,text width=\dimexpr0.200000\linewidth-2.5mm\relax] at (0.900000,3.250) {$69$及$69$以下};
\node[font=\normalsize,align=center,inner sep=0pt,text width=\dimexpr0.200000\linewidth-2.5mm\relax] at (0.100000,9.750) {四（1）班};
\node[font=\normalsize,align=center,inner sep=0pt,text width=\dimexpr0.200000\linewidth-2.5mm\relax] at (0.300000,9.750) {{\color{answerred}$2$}};
\node[font=\normalsize,align=center,inner sep=0pt,text width=\dimexpr0.200000\linewidth-2.5mm\relax] at (0.500000,9.750) {{\color{answerred}$5$}};
\node[font=\normalsize,align=center,inner sep=0pt,text width=\dimexpr0.200000\linewidth-2.5mm\relax] at (0.700000,9.750) {{\color{answerred}$2$}};
\node[font=\normalsize,align=center,inner sep=0pt,text width=\dimexpr0.200000\linewidth-2.5mm\relax] at (0.900000,9.750) {{\color{answerred}$1$}};
\node[font=\normalsize,align=center,inner sep=0pt,text width=\dimexpr0.200000\linewidth-2.5mm\relax] at (0.100000,16.250) {四（2）班};
\node[font=\normalsize,align=center,inner sep=0pt,text width=\dimexpr0.200000\linewidth-2.5mm\relax] at (0.300000,16.250) {{\color{answerred}$4$}};
\node[font=\normalsize,align=center,inner sep=0pt,text width=\dimexpr0.200000\linewidth-2.5mm\relax] at (0.500000,16.250) {{\color{answerred}$3$}};
\node[font=\normalsize,align=center,inner sep=0pt,text width=\dimexpr0.200000\linewidth-2.5mm\relax] at (0.700000,16.250) {{\color{answerred}$2$}};
\node[font=\normalsize,align=center,inner sep=0pt,text width=\dimexpr0.200000\linewidth-2.5mm\relax] at (0.900000,16.250) {{\color{answerred}$1$}};
\end{tikzpicture}\par\endgroup

\QHeading{分类核对}

\par\noindent\begingroup
\begin{tikzpicture}[x=\dimexpr\linewidth-1pt\relax,y=-1mm,draw=black,text=black,line width=.45pt]
\draw (0.000000,0)--(0.000000,43.000);
\draw (0.200000,0)--(0.200000,43.000);
\draw (0.600000,0)--(0.600000,43.000);
\draw (1.000000,0)--(1.000000,43.000);
\draw (0,0.000)--(1,0.000);
\draw (0,7.000)--(1,7.000);
\draw (0,17.000)--(1,17.000);
\draw (0,27.000)--(1,27.000);
\draw (0,35.000)--(1,35.000);
\draw (0,43.000)--(1,43.000);
\node[font=\normalsize,align=center,inner sep=0pt,text width=\dimexpr0.200000\linewidth-2.5mm\relax] at (0.100000,3.500) {成绩/分};
\node[font=\normalsize,align=center,inner sep=0pt,text width=\dimexpr0.400000\linewidth-2.5mm\relax] at (0.400000,3.500) {四（1）班};
\node[font=\normalsize,align=center,inner sep=0pt,text width=\dimexpr0.400000\linewidth-2.5mm\relax] at (0.800000,3.500) {四（2）班};
\node[font=\normalsize,align=center,inner sep=0pt,text width=\dimexpr0.200000\linewidth-2.5mm\relax] at (0.100000,12.000) {$90$及$90$以上};
\node[font=\normalsize,align=center,inner sep=0pt,text width=\dimexpr0.400000\linewidth-2.5mm\relax] at (0.400000,12.000) {{\color{answerred}$92$、$95$（$2$人）}};
\node[font=\normalsize,align=center,inner sep=0pt,text width=\dimexpr0.400000\linewidth-2.5mm\relax] at (0.800000,12.000) {{\color{answerred}$95$、$98$、$90$、$92$（$4$人）}};
\node[font=\normalsize,align=center,inner sep=0pt,text width=\dimexpr0.200000\linewidth-2.5mm\relax] at (0.100000,22.000) {$80\sim89$};
\node[font=\normalsize,align=center,inner sep=0pt,text width=\dimexpr0.400000\linewidth-2.5mm\relax] at (0.400000,22.000) {{\color{answerred}$85$、$87$、$87$、$85$、$89$（$5$人）}};
\node[font=\normalsize,align=center,inner sep=0pt,text width=\dimexpr0.400000\linewidth-2.5mm\relax] at (0.800000,22.000) {{\color{answerred}$87$、$85$、$87$（$3$人）}};
\node[font=\normalsize,align=center,inner sep=0pt,text width=\dimexpr0.200000\linewidth-2.5mm\relax] at (0.100000,31.000) {$70\sim79$};
\node[font=\normalsize,align=center,inner sep=0pt,text width=\dimexpr0.400000\linewidth-2.5mm\relax] at (0.400000,31.000) {{\color{answerred}$78$、$72$（$2$人）}};
\node[font=\normalsize,align=center,inner sep=0pt,text width=\dimexpr0.400000\linewidth-2.5mm\relax] at (0.800000,31.000) {{\color{answerred}$73$、$77$（$2$人）}};
\node[font=\normalsize,align=center,inner sep=0pt,text width=\dimexpr0.200000\linewidth-2.5mm\relax] at (0.100000,39.000) {$69$及$69$以下};
\node[font=\normalsize,align=center,inner sep=0pt,text width=\dimexpr0.400000\linewidth-2.5mm\relax] at (0.400000,39.000) {{\color{answerred}$60$（$1$人）}};
\node[font=\normalsize,align=center,inner sep=0pt,text width=\dimexpr0.400000\linewidth-2.5mm\relax] at (0.800000,39.000) {{\color{answerred}$66$（$1$人）}};
\end{tikzpicture}\par\endgroup
\Answer{四（2）班的$90$分也属于“$90$及$90$以上”，该格应填$4$，不是照片手写的$3$。两个班各$10$人。}
\Problem{② 分别求出这两个班参加知识竞赛的同学的平均成绩。}
\Answer{四（1）班（运用加法交换律、结合律分组计算）：}
\Answer{\Formula{\begin{aligned}[t]&\phantom{={}}(85+92+87+87+85+60+95+89+78+72)\div10\\
&=[(85+95)+(92+78)+(87+87)+(85+89)+(60+72)]\div10\\
&=[180+170+174+174+132]\div10\\
&=[(180+170)+(174+174)+132]\div10\\
&=[350+348+132]\div10\\
&=[350+(348+132)]\div10\\
&=[350+480]\div10\\
&=830\div10\\
&=83\text{（分）}\end{aligned}}}
\Answer{答：四（1）班参加知识竞赛的同学平均成绩是$83$分。}
\Answer{四（2）班（运用加法交换律、结合律分组计算）：}
\Answer{\Formula{\begin{aligned}[t]&\phantom{={}}(66+95+87+98+73+77+85+87+90+92)\div10\\
&=[(95+85)+(87+73)+(98+92)+(66+90)+(77+87)]\div10\\
&=[180+160+190+156+164]\div10\\
&=[180+160+190+(156+164)]\div10\\
&=[180+160+190+320]\div10\\
&=[(180+320)+(160+190)]\div10\\
&=[500+350]\div10\\
&=850\div10\\
&=85\text{（分）}\end{aligned}}}
\Answer{答：四（2）班参加知识竞赛的同学平均成绩是$85$分。}

\QHeading{第2题②　方法二：以80分为基数}

\par\noindent\begin{minipage}[t]{.485\linewidth}
\Answer{四（1）班多出量合计：}
\Answer{\Formula{\begin{aligned}[t]&\phantom{={}}5+12+7+7+5+15+9\\
&=(5+15)+(12+7+7)+(5+9)\\
&=20+(19+7)+14\\
&=20+26+14\\
&=20+(26+14)\\
&=20+40\\
&=60\text{（分）}\end{aligned}}}
\Answer{不足量合计：}
\Answer{\Formula{\begin{aligned}[t]&\phantom{={}}20+2+8\\
&=20+(2+8)\\
&=20+10\\
&=30\text{（分）}\end{aligned}}}
\Answer{平均成绩：}
\Answer{\Formula{\begin{aligned}[t]&\phantom{={}}80+(60-30)\div10\\
&=80+30\div10\\
&=80+3\\
&=83\text{（分）}\end{aligned}}}
\Answer{答：四（1）班平均成绩是$83$分。}

\end{minipage}\hfill\begin{minipage}[t]{.485\linewidth}
\Answer{四（2）班多出量合计：}
\Answer{\Formula{\begin{aligned}[t]&\phantom{={}}15+7+18+5+7+10+12\\
&=(15+5)+(18+12)+(7+7)+10\\
&=20+30+14+10\\
&=(20+30)+(14+10)\\
&=50+24\\
&=74\text{（分）}\end{aligned}}}
\Answer{不足量合计：}
\Answer{\Formula{\begin{aligned}[t]&\phantom{={}}14+7+3\\
&=14+(7+3)\\
&=14+10\\
&=24\text{（分）}\end{aligned}}}
\Answer{平均成绩：}
\Answer{\Formula{\begin{aligned}[t]&\phantom{={}}80+(74-24)\div10\\
&=80+50\div10\\
&=80+5\\
&=85\text{（分）}\end{aligned}}}
\Answer{答：四（2）班平均成绩是$85$分。}

\end{minipage}\par
